\documentclass[11pt, a4paper]{article}

% bfw-style.tex — gemeinsame Style-Definitionen für alle Handouts/Aufgabenhefte
% Voraussetzung: Kompilation mit xelatex (wegen fontspec)

\usepackage[ngerman]{babel}
\usepackage{geometry}
\geometry{a4paper, top=2.2cm, bottom=2.2cm, left=2.3cm, right=2.3cm}

\usepackage{fontspec}
% Fallback-Kette: Source Sans Pro -> Calibri -> Segoe UI -> Latin Modern Sans
% (Latin Modern ist immer mit MiKTeX vorhanden)
\IfFontExistsTF{Source Sans Pro}{%
  \setmainfont{Source Sans Pro}[Ligatures=TeX]%
}{\IfFontExistsTF{Calibri}{%
  \setmainfont{Calibri}[Ligatures=TeX]%
}{\IfFontExistsTF{Segoe UI}{%
  \setmainfont{Segoe UI}[Ligatures=TeX]%
}{%
  \setmainfont{Latin Modern Sans}[Ligatures=TeX]%
}}}
\IfFontExistsTF{Source Code Pro}{%
  \setmonofont{Source Code Pro}[Scale=0.92]%
}{\IfFontExistsTF{Consolas}{%
  \setmonofont{Consolas}[Scale=0.92]%
}{%
  \setmonofont{Latin Modern Mono}[Scale=0.92]%
}}

\usepackage{xcolor}
% Farbpalette: hell, freundlich, modern
\definecolor{bfwPrimary}{HTML}{0EA5A4}    % Petrol/Türkis - Primär
\definecolor{bfwPrimaryDark}{HTML}{0F766E} % dunkler Petrol für Headings
\definecolor{bfwAccent}{HTML}{F59E0B}     % Amber - Akzent
\definecolor{bfwSuccess}{HTML}{10B981}    % Grün - Erfolg / Lösung
\definecolor{bfwWarn}{HTML}{F97316}       % Orange - Achtung
\definecolor{bfwInk}{HTML}{0F172A}        % fast Schwarz
\definecolor{bfwInkSoft}{HTML}{475569}    % gedämpftes Grau
\definecolor{bfwBgSoft}{HTML}{F8FAFC}     % off-white
\definecolor{bfwBgTeal}{HTML}{ECFEFF}     % helles Türkis
\definecolor{bfwBgAmber}{HTML}{FEF3C7}    % helles Gelb
\definecolor{bfwBgRose}{HTML}{FFE4E6}     % helles Rosé für Eselsbrücken

\usepackage[colorlinks=true, linkcolor=bfwPrimaryDark, urlcolor=bfwPrimaryDark]{hyperref}
\usepackage{enumitem}
\setlist[itemize]{leftmargin=*, itemsep=2pt, topsep=2pt}
\setlist[enumerate]{leftmargin=*, itemsep=2pt, topsep=2pt}

\usepackage[most]{tcolorbox}
\usepackage{booktabs}
\usepackage{tikz}
\usetikzlibrary{decorations.pathmorphing, positioning, shapes.geometric, arrows.meta}

% Merkkästchen — wichtige Definition / Kernpunkt
\newtcolorbox{merksatz}[1][]{%
  enhanced, breakable,
  colback=bfwBgTeal,
  colframe=bfwPrimary,
  boxrule=0.6pt,
  arc=4pt,
  left=10pt, right=10pt, top=8pt, bottom=8pt,
  fonttitle=\bfseries\color{bfwPrimaryDark},
  title={\faLightbulb\ Merksatz},
  #1
}

% Eselsbrücke — Mnemoniks
\newtcolorbox{eselsbruecke}[1][]{%
  enhanced, breakable,
  colback=bfwBgRose,
  colframe=bfwAccent,
  boxrule=0.6pt,
  arc=4pt,
  left=10pt, right=10pt, top=8pt, bottom=8pt,
  fonttitle=\bfseries\color{bfwWarn},
  title={Eselsbrücke},
  #1
}

% Beispiel-Box
\newtcolorbox{beispiel}[1][]{%
  enhanced, breakable,
  colback=bfwBgSoft,
  colframe=bfwInkSoft,
  boxrule=0.4pt,
  arc=3pt,
  left=10pt, right=10pt, top=8pt, bottom=8pt,
  fonttitle=\bfseries\color{bfwInk},
  title={Beispiel},
  #1
}

% Lösungs-Box (für Aufgabenhefte)
\newtcolorbox{loesung}[1][]{%
  enhanced, breakable,
  colback=bfwBgAmber!40,
  colframe=bfwSuccess,
  boxrule=0.6pt,
  arc=3pt,
  left=10pt, right=10pt, top=8pt, bottom=8pt,
  fonttitle=\bfseries\color{bfwSuccess!70!black},
  title={Lösung},
  #1
}

% Glossar-Eintrag
\newcommand{\glossareintrag}[2]{%
  \par\noindent\textbf{\color{bfwPrimaryDark}#1} \enspace #2\par\smallskip
}

% Section-Stil — etwas farbiger als Standard
\usepackage{titlesec}
\titleformat{\section}
  {\Large\bfseries\color{bfwPrimaryDark}}
  {\thesection}{0.6em}{}
\titleformat{\subsection}
  {\large\bfseries\color{bfwPrimary}}
  {\thesubsection}{0.5em}{}
\titleformat{\subsubsection}
  {\normalsize\bfseries\color{bfwInk}}
  {\thesubsubsection}{0.4em}{}

% TikZ-Stil "skizzenhaft" — etwas weicher als Rough.js, aber stabil
\tikzset{
  roughbox/.style = {
    draw=bfwPrimaryDark, thick, fill=bfwBgTeal,
    rounded corners=4pt, inner sep=6pt
  },
  roughaccent/.style = {
    draw=bfwAccent, thick, fill=bfwBgAmber,
    rounded corners=4pt, inner sep=6pt
  },
  roughline/.style = {
    draw=bfwInkSoft, thick, -{Stealth[length=2.5mm]}
  }
}

% Bullet-Symbole (bewusst kein FontAwesome — vermeidet Font-Installations-Probleme)
\newcommand{\faLightbulb}{$\bullet$}
\newcommand{\faBookmark}{$\bullet$}

% Header/Footer
\usepackage{fancyhdr}
\pagestyle{fancy}
\fancyhf{}
\renewcommand{\headrulewidth}{0pt}
\fancyfoot[L]{\small\color{bfwInkSoft}\thepage}
\fancyfoot[R]{\small\color{bfwInkSoft} BFW M-064 Beschaffung im Büromanagement}


\title{\vspace*{-1.5cm}\Huge\bfseries\color{bfwPrimaryDark}Tag~7: Eingangsrechnungen\\[0.4em]
  \large\color{bfwInkSoft}\textnormal{Erfüllung, Eigentumsvorbehalt, Gefahrübergang, Rechnungsprüfung}}
\author{\textsf{Beschaffung im Büromanagement}\\
\textsf{\small\copyright\ Can Siebert\,\textperiodcentered\,2026}}
\date{20.05.2026}

\begin{document}
\maketitle
\thispagestyle{empty}

\vspace{1em}
\begin{merksatz}[title={\bfseries Worum geht es heute?}]
Die Ware ist da, jetzt kommt die Rechnung. Bevor die Buchhaltung zahlt, muss die
Eingangsrechnung sorgfältig geprüft werden --- formell, sachlich, rechnerisch.
Daneben klären wir wichtige rechtliche Fragen: Wann ist der Kaufvertrag erfüllt?
Wem gehört die Ware? Wer trägt die Gefahr beim Versand? Genau diese Themen werden
in der IHK-Prüfung Teil~1 (Bürokommunikation, Tabellenkalkulation) und Teil~2
(Wirtschafts- und Sozialkunde) regelmäßig abgefragt.
\end{merksatz}

\tableofcontents
\vspace{1em}
\hrule
\vspace{1em}

\section{Erfüllung des Kaufvertrags}

\subsection{§ 362 BGB --- Erlöschen durch Leistung}

Der Kaufvertrag wird mit der \emph{Erfüllung} der geschuldeten Leistungen beendet.
Nach §\,362 Abs.\,1 BGB erlischt das Schuldverhältnis, wenn die geschuldete Leistung
\emph{an den Gläubiger} bewirkt wird.

Beim Kaufvertrag bedeutet das:

\begin{itemize}[itemsep=2pt]
  \item Der Verkäufer erfüllt durch \emph{Übergabe der Sache} und \emph{Verschaffung des Eigentums} (§\,433 I BGB).
  \item Der Käufer erfüllt durch \emph{Zahlung des Kaufpreises} und \emph{Abnahme der Sache} (§\,433 II BGB).
\end{itemize}

\section{Gutgläubiger Eigentumserwerb --- § 932 BGB}

\subsection{Was ist gutgläubiger Erwerb?}

Normalerweise kann nur der Eigentümer das Eigentum übertragen. § 932 BGB durchbricht
diese Regel: Wer eine bewegliche Sache von einem \emph{Nichtberechtigten} erwirbt,
kann gleichwohl Eigentümer werden, wenn er bei der Übergabe \emph{guten Glaubens} ist.

\subsection{Voraussetzungen}

\begin{enumerate}[itemsep=2pt]
  \item Verkehrsgeschäft (Erwerb von einem anderen).
  \item Bewegliche Sache.
  \item Übergabe (oder Übergabesurrogat).
  \item Guter Glaube an das Eigentum des Veräußerers (§ 932 II BGB --- nicht grob fahrlässige Unkenntnis).
\end{enumerate}

\subsection{Ausnahme: Abhandengekommene Sachen}

\begin{merksatz}[title={§ 935 BGB}]
Bei \emph{gestohlenen, verlorenen oder sonst abhandengekommenen} Sachen ist
gutgläubiger Erwerb \emph{ausgeschlossen} --- es sei denn, die Sache wird in einer
öffentlichen Versteigerung erworben.
\end{merksatz}

\section{Eigentumsvorbehalt --- § 449 BGB}

\subsection{Idee und Grundform}

Der \emph{Eigentumsvorbehalt} schützt den Verkäufer: Das Eigentum geht erst dann
auf den Käufer über, wenn der Kaufpreis vollständig bezahlt ist. Bis dahin bleibt
der Verkäufer Eigentümer, der Käufer ist ,,nur'' Besitzer.

\subsection{Vier Formen}

\begin{center}
\begin{tabular}{p{4cm} p{10cm}}
\toprule
\textbf{Form} & \textbf{Inhalt}\\
\midrule
Einfacher EV & Eigentum geht erst mit voller Bezahlung der konkreten Lieferung über.\\
Erweiterter EV & Eigentum geht erst über, wenn \emph{alle} Forderungen aus der Geschäftsverbindung beglichen sind (Kontokorrent-Vorbehalt).\\
Verlängerter EV & Käufer darf weiterverkaufen; im Voraus tritt er die Forderungen aus dem Weiterverkauf an den Lieferanten ab (§ 398 BGB).\\
Weitergeleiteter EV & Vorbehaltsrecht wird auf einen Dritten übertragen (z.\,B.\ bei Verarbeitung in eine Maschine eines anderen Eigentümers).\\
\bottomrule
\end{tabular}
\end{center}

\begin{eselsbruecke}
,,\textbf{EVE-W}'': \textbf{E}infach --- \textbf{V}erlängert --- \textbf{E}rweitert --- \textbf{W}eitergeleitet. So merken Sie sich die vier Formen.
\end{eselsbruecke}

\section{Erfüllungsort, Gefahrübergang, Gerichtsstand}

\subsection{Erfüllungsort --- § 269 BGB}

Der \emph{Erfüllungsort} bestimmt, \emph{wo} der Schuldner die Leistung zu erbringen hat.
Drei Konstellationen:

\begin{itemize}[itemsep=2pt]
  \item \textbf{Holschuld}: Erfüllungsort ist beim \emph{Schuldner} (Verkäufer). Käufer holt die Ware ab. \emph{Gesetzlicher Regelfall.}
  \item \textbf{Bringschuld}: Erfüllungsort ist beim \emph{Gläubiger} (Käufer). Verkäufer bringt die Ware hin.
  \item \textbf{Schickschuld}: Erfüllungsort beim Schuldner, aber er versendet auf Käufergefahr (Versendungskauf).
\end{itemize}

\subsection{Gefahrübergang --- §§ 446, 447 BGB}

Die \emph{Preisgefahr} fragt: Wer trägt das Risiko, wenn die Sache zufällig untergeht
(z.\,B.\ beim Transport zerstört wird)?

\begin{merksatz}[title={§ 446 BGB --- Grundregel}]
Mit \emph{Übergabe} der Kaufsache geht die Gefahr auf den Käufer über. Bis zur
Übergabe trägt der Verkäufer das Risiko.
\end{merksatz}

\begin{merksatz}[title={§ 447 BGB --- Versendungskauf}]
Beim Versendungskauf (Versand auf Verlangen des Käufers an einen anderen Ort) geht
die Gefahr bereits mit \emph{Übergabe an den Spediteur/Frachtführer} über.
\end{merksatz}

\textbf{Achtung Verbraucher:} Bei einem Verbrauchsgüterkauf gilt § 475 II BGB --- die
Gefahr geht erst mit Übergabe an den Verbraucher über, auch wenn versendet wird.

\subsection{Gerichtsstand}

Bei kaufmännischen Streitigkeiten ist häufig ein \emph{Gerichtsstand} im Vertrag
vereinbart (zulässig zwischen Kaufleuten, § 38 ZPO). Sonst gilt der allgemeine
Gerichtsstand: Wohnsitz des Beklagten (§§ 12, 13 ZPO).

\section{Pflichtangaben einer Rechnung --- § 14 UStG}

Eine Rechnung muss alle Pflichtangaben enthalten, damit der Empfänger den
\emph{Vorsteuerabzug} geltend machen kann.

\begin{merksatz}[title={Die wichtigsten Pflichtangaben}]
\begin{enumerate}[itemsep=0pt]
  \item Vollständiger Name und Anschrift von Lieferant und Empfänger
  \item Steuernummer oder USt-Identifikationsnummer des Lieferanten
  \item Ausstellungsdatum der Rechnung
  \item Fortlaufende, einmalige Rechnungsnummer
  \item Menge und handelsübliche Bezeichnung der Lieferung
  \item Liefer- oder Leistungsdatum
  \item Nettoentgelt nach Steuersätzen aufgeschlüsselt
  \item Anzuwendender Steuersatz und Steuerbetrag (oder Hinweis auf Steuerbefreiung)
\end{enumerate}
\end{merksatz}

Bei \emph{Kleinbetragsrechnungen} bis 250\,€ brutto reichen vereinfachte Angaben
(§ 33 UStDV).

\section{Rechnungsprüfung in der Praxis}

\subsection{Drei Stufen der Prüfung}

\begin{center}
\begin{tabular}{p{3.5cm} p{10.5cm}}
\toprule
\textbf{Stufe} & \textbf{Inhalt}\\
\midrule
Formell & Sind alle Pflichtangaben nach § 14 UStG vorhanden?\\
Sachlich & Stimmt die Rechnung mit Bestellung und Wareneingang überein? Wurde die Leistung tatsächlich erbracht?\\
Rechnerisch & Sind die Beträge richtig? Mengen $\times$ Preise, Zwischensummen, USt, Brutto, Skonto/Rabatt?\\
\bottomrule
\end{tabular}
\end{center}

\subsection{Drei-Punkt-Vergleich}

\begin{merksatz}[title={Drei-Punkt-Vergleich}]
\textbf{Bestellung} (Was war vereinbart?) $\to$ \textbf{Wareneingang/Lieferschein} (Was ist gekommen?) $\to$ \textbf{Rechnung} (Was steht in der Rechnung?). Alle drei müssen übereinstimmen.
\end{merksatz}

\subsection{Beispielrechnung}

\begin{beispiel}
Bestellung: 50 Aktenordner à 4,80\,€, 10 Pakete Druckerpapier à 28,00\,€, Versand 12\,€, USt 19\,\%.

\begin{itemize}[itemsep=0pt]
  \item Aktenordner: $50 \times 4{,}80 = \mathbf{240{,}00~€}$
  \item Druckerpapier: $10 \times 28{,}00 = \mathbf{280{,}00~€}$
  \item Versand: $\mathbf{12{,}00~€}$
  \item Zwischensumme netto: $\mathbf{532{,}00~€}$
  \item USt 19\,\%: $\mathbf{101{,}08~€}$
  \item Bruttosumme: $\mathbf{633{,}08~€}$
\end{itemize}
\end{beispiel}

\subsection{Was tun bei Fehlern?}

\begin{itemize}[itemsep=2pt]
  \item Formelle Fehler $\to$ Rechnung an den Lieferanten zur Korrektur zurücksenden (Vorsteuerabzug nicht möglich).
  \item Sachliche Fehler $\to$ klären (z.\,B.\ falsche Menge), ggf.\ Nachbelastung oder Gutschrift.
  \item Rechnerische Fehler $\to$ korrigierte Rechnung anfordern.
\end{itemize}

\section{Zusammenfassung}

\begin{enumerate}[itemsep=2pt]
  \item \textbf{Erfüllung} (§ 362 BGB): Vertrag erlischt mit ordnungsgemäßer Leistung.
  \item \textbf{Gutgläubiger Erwerb} (§ 932 BGB): Eigentum vom Nichtberechtigten möglich --- nicht bei abhandengekommenen Sachen (§ 935).
  \item \textbf{Eigentumsvorbehalt}: einfach, erweitert, verlängert, weitergeleitet.
  \item \textbf{Gefahrübergang}: § 446 mit Übergabe; § 447 mit Übergabe an Spediteur (B2B); § 475 II für Verbraucher anders.
  \item \textbf{Pflichtangaben} § 14 UStG --- ohne sie kein Vorsteuerabzug.
  \item \textbf{Drei-Stufen-Prüfung}: formell, sachlich, rechnerisch.
\end{enumerate}

\newpage

\section*{Glossar}
\addcontentsline{toc}{section}{Glossar}

\glossareintrag{Erfüllung}{Erlöschen des Schuldverhältnisses durch Bewirkung der geschuldeten Leistung (§ 362 BGB).}

\glossareintrag{Gutgläubiger Erwerb}{Eigentumserwerb vom Nichtberechtigten, wenn der Erwerber gutgläubig ist (§ 932 BGB).}

\glossareintrag{§ 935 BGB}{Ausschluss des gutgläubigen Erwerbs bei abhandengekommenen Sachen (Diebstahl, Verlust).}

\glossareintrag{Eigentumsvorbehalt}{Klausel im Kaufvertrag (§ 449 BGB): Eigentum geht erst mit voller Bezahlung über.}

\glossareintrag{Einfacher EV}{Eigentum geht mit Zahlung der konkreten Lieferung über.}

\glossareintrag{Erweiterter EV}{Eigentum geht erst mit Begleichung aller offenen Forderungen aus der Geschäftsverbindung über.}

\glossareintrag{Verlängerter EV}{Käufer darf weiterverkaufen; Forderungen werden im Voraus an den Lieferanten abgetreten (§ 398 BGB).}

\glossareintrag{Weitergeleiteter EV}{Vorbehaltsrecht wird an einen Dritten übertragen (z.\,B.\ bei Verarbeitung der Ware).}

\glossareintrag{Erfüllungsort}{Ort, an dem der Schuldner die Leistung zu erbringen hat (§ 269 BGB).}

\glossareintrag{Holschuld}{Erfüllungsort beim Schuldner --- Käufer holt ab. Gesetzlicher Regelfall.}

\glossareintrag{Bringschuld}{Erfüllungsort beim Gläubiger --- Verkäufer bringt die Ware hin.}

\glossareintrag{Schickschuld / Versendungskauf}{Erfüllungsort beim Schuldner, aber Versand auf Käufergefahr (§ 447 BGB).}

\glossareintrag{Gefahrübergang}{Übergang der Preisgefahr vom Verkäufer auf den Käufer (§ 446 BGB Grundregel; § 447 BGB Versendungskauf).}

\glossareintrag{§ 14 UStG}{Vorschrift zu den Pflichtangaben einer Rechnung --- Voraussetzung für den Vorsteuerabzug.}

\glossareintrag{Vorsteuerabzug}{Recht des Unternehmers, gezahlte USt aus Eingangsrechnungen mit der eigenen USt-Schuld zu verrechnen.}

\glossareintrag{Steuernummer / USt-IdNr.}{Pflichtangabe nach § 14 IV Nr. 2 UStG --- ohne diese kein Vorsteuerabzug.}

\glossareintrag{Formelle Prüfung}{Prüfung, ob alle Pflichtangaben nach § 14 UStG vollständig auf der Rechnung sind.}

\glossareintrag{Sachliche Prüfung}{Prüfung, ob die Rechnung mit Bestellung und Lieferung übereinstimmt.}

\glossareintrag{Rechnerische Prüfung}{Prüfung, ob Beträge, Zwischensummen, USt und Bruttosummen korrekt berechnet sind.}

\glossareintrag{Drei-Punkt-Vergleich}{Abgleich von Bestellung, Wareneingang/Lieferschein und Rechnung in der Rechnungsprüfung.}

\glossareintrag{Kleinbetragsrechnung}{Rechnung bis 250\,€ brutto --- vereinfachte Pflichtangaben nach § 33 UStDV.}

\end{document}
